\begin{exercises}
  \item
    转置以下各矩阵。
    \begin{exparts*}
      \partsitem \( \begin{mat}[r]
                 2  &1  \\
                 3  &1
               \end{mat}  \)
      \partsitem \( \begin{mat}[r]
                 2  &1  \\
                 1  &3
               \end{mat}  \)
      \partsitem \( \begin{mat}[r]
                 1  &4  &3 \\
                 6  &7  &8
               \end{mat}  \)
      \partsitem \( \colvec[r]{0 \\ 0 \\ 0} \)
      \partsitem \( \rowvec{-1 &-2} \)
    \end{exparts*}
    \begin{answer}
      \begin{exparts*}
        \partsitem \( \begin{mat}[r]
                   2  &3  \\
                   1  &1
                 \end{mat}  \)
        \partsitem \( \begin{mat}[r]
                   2  &1  \\
                   1  &3
                 \end{mat}  \)
        \partsitem \( \begin{mat}[r]
                   1  &6  \\
                   4  &7  \\
                   3  &8
                 \end{mat}  \)
        \partsitem \( \rowvec{0 &0 &0} \)
        \partsitem \( \colvec[r]{-1 \\ -2} \)
      \end{exparts*}
    \end{answer}
  \recommended \item
    判断该向量是否在矩阵的行空间中。
    \begin{exparts*}
       \partsitem \( \begin{mat}[r]
                  2  &1  \\
                  3  &1
                \end{mat}  \),
             \( \rowvec{1 &0} \)
       \partsitem \( \begin{mat}[r]
                  0  &1  &3  \\
                 -1  &0  &1  \\
                 -1  &2  &7
                \end{mat}  \),
             \( \rowvec{1 &1 &1} \)
    \end{exparts*}
    \begin{answer}
       \begin{exparts}
         \partsitem 是的。
           为判断是否存在 $c_1$ 与 $c_2$ 使得
           \( c_1\cdot\rowvec{2 &1}+c_2\cdot\rowvec{3 &1}=\rowvec{1 &0} \)，
           我们求解
           \begin{equation*}
              \begin{linsys}{2}
                 2c_1  &+  &3c_2  &=  &1  \\
                  c_1  &+  &c_2   &=  &0
              \end{linsys}
           \end{equation*}
           解得 \( c_1=-1 \) 与 \( c_2=1 \)。
           因此该向量在行空间中。
         \partsitem 否。
           方程
           \( c_1\rowvec{0 &1 &3}
              +c_2\rowvec{-1 &0 &1}
              +c_3\rowvec{-1 &2 &7}
              =\rowvec{1 &1 &1} \)
           无解。
           \begin{equation*}
             \begin{amat}[r]{3}
               0  &-1  &-1  &1  \\
               1  &0   &2   &1  \\
               3  &1   &7   &1
             \end{amat}
             \grstep{\rho_1\leftrightarrow\rho_2}
             \grstep{-3\rho_1+\rho_3}
             \repeatedgrstep{\rho_2+\rho_3}
             \begin{amat}[r]{3}
               1  &0   &2   &1  \\
               0  &-1  &-1  &1  \\
               0  &0   &0   &-1
             \end{amat}
           \end{equation*}
           因此，该向量不在行空间中。
      \end{exparts}
    \end{answer}
  \recommended \item
    判断该向量是否在列空间中。
    \begin{exparts*}
      \partsitem \( \begin{mat}[r]
                 1  &1  \\
                 1  &1
               \end{mat}  \),
            \( \colvec[r]{1 \\ 3} \)
      \partsitem \( \begin{mat}[r]
                 1  &3  &1 \\
                 2  &0  &4 \\
                 1  &-3 &-3
               \end{mat}  \),
            \( \colvec[r]{1 \\ 0 \\ 0} \)
    \end{exparts*}
    \begin{answer}
       \begin{exparts}
          \partsitem 否。
            为判断是否存在 \( c_1,c_2\in\Re \) 使得
            \begin{equation*}
              c_1\colvec[r]{1 \\ 1}
              +c_2\colvec[r]{1 \\ 1}
              =\colvec[r]{1 \\ -3}
            \end{equation*}
            我们可以对由此得到的线性方程组使用高斯消元法。
            \begin{equation*}
              \begin{linsys}{2}
                 c_1  &+  &c_2  &=  &1  \\
                 c_1  &+  &c_2  &=  &3
              \end{linsys}
              \grstep{-\rho_1+\rho_2}
              \begin{linsys}{2}
                 c_1  &+  &c_2  &=  &1  \\
                      &   &0    &=  &2
              \end{linsys}
            \end{equation*}
            该方程组无解，因此该向量不在列空间中。
          \partsitem 是的。
            由关系式
            \begin{equation*}
              c_1\colvec[r]{1 \\ 2 \\ 1}
              +c_2\colvec[r]{3 \\ 0 \\ -3}
              +c_3\colvec[r]{1 \\ 4 \\ -3}
              =\colvec[r]{1 \\ 0 \\ 0}
            \end{equation*}
            我们得到一个线性方程组，对其使用高斯消元法，
            \begin{equation*}
              \begin{amat}[r]{3}
                1  &3  &1  &1  \\
                2  &0  &4  &0  \\
                1  &-3 &-3 &0
              \end{amat}
              \grstep[-\rho_1+\rho_3]{-2\rho_1+\rho_2}
              \repeatedgrstep{-\rho_2+\rho_3}
              \begin{amat}[r]{3}
                1  &3  &1  &1  \\
                0  &-6 &2  &-2 \\
                0  &0  &-6 &1
              \end{amat}
            \end{equation*}
            即得一个解。
            因此，该向量在列空间中。
% sage: load('gauss_method.sage')
% sage: M = matrix(QQ, [[1,3, 1, 1], [2, 0, 4, 0], [1, -3, -3, 0]])
% sage: M
% [ 1  3  1  1]
% [ 2  0  4  0]
% [ 1 -3 -3  0]
% sage: gauss_method(M)
% [ 1  3  1  1]
% [ 2  0  4  0]
% [ 1 -3 -3  0]
% (' take', -2, 'times row', 1, 'plus row', 2)
% (' take', -1, 'times row', 1, 'plus row', 3)
% [ 1  3  1  1]
% [ 0 -6  2 -2]
% [ 0 -6 -4 -1]
% (' take', -1, 'times row', 2, 'plus row', 3)
% [ 1  3  1  1]
% [ 0 -6  2 -2]
% [ 0  0 -6  1]
      \end{exparts}
    \end{answer}
  \recommended \item
    判断该向量是否在矩阵的列空间中。
    \begin{exparts*}
      \partsitem \( \begin{mat}[r]
                   2  &1  \\
                   2  &5
                 \end{mat} \),~\( \colvec[r]{1 \\ -3}  \)
      \partsitem \( \begin{mat}[r]
                   4  &-8 \\
                   2  &-4
                 \end{mat} \),~\( \colvec[r]{0 \\ 1}  \)
      \partsitem \( \begin{mat}[r]
                   1  &-1  &1  \\
                   1  &1   &-1 \\
                  -1  &-1  &1
                \end{mat} \),~\( \colvec[r]{2 \\ 0 \\ 0}  \)
    \end{exparts*}
    \begin{answer}
      \begin{exparts}
        \partsitem 是的；
          我们问的是是否存在标量 \( c_1 \) 与 \( c_2 \) 使得
          \begin{equation*}
            c_1\colvec[r]{2 \\ 2}+c_2\colvec[r]{1 \\ 5}=\colvec[r]{1 \\ -3}
          \end{equation*}
          由此得到线性方程组
          \begin{equation*}
            \begin{linsys}{2}
              2c_1  &+  &c_2  &=  &1  \\
              2c_1  &+  &5c_2 &=  &-3
            \end{linsys}
            \grstep{-\rho_1+\rho_2}
            \begin{linsys}{2}
              2c_1  &+  &c_2  &=  &1  \\
                    &   &4c_2 &=  &-4
            \end{linsys}
          \end{equation*}
          高斯消元法给出 \( c_2=-1 \) 与 \( c_1=1 \)。
          也就是说，确实存在这样一对标量，因此该向量
          确实在矩阵的列空间中。
        \partsitem 否；
          我们问的是是否存在标量 $c_1$ 与 $c_2$
          使得
          \begin{equation*}
            c_1\colvec[r]{4 \\ 2}+c_2\colvec[r]{-8 \\ -4}=\colvec[r]{0 \\ 1}
          \end{equation*}
          一种做法是考虑由此得到的线性方程组
          \begin{equation*}
            \begin{linsys}{2}
              4c_1  &-  &8c_2  &=  &0  \\
              2c_1  &-  &4c_2  &=  &1
            \end{linsys}
          \end{equation*}
          容易看出它无解。
          另一种做法是注意到
          左边各列的任意线性组合
          其第二个分量都是第一个分量的一半，
          而右边的向量不满足这一条件。
        \partsitem 是的；我们可以直接观察到该向量
          是第一列减去第二列。
          或者，若一时看不出，可以写出各列之间的关系
          \begin{equation*}
            c_1\colvec[r]{1 \\ 1 \\ -1}
             +c_2\colvec[r]{-1 \\ 1 \\ -1}
             +c_3\colvec[r]{1 \\ -1 \\ 1}
             =\colvec[r]{2 \\ 0 \\ 0}
          \end{equation*}
          并考虑由此得到的线性方程组
          \begin{equation*}
            \begin{linsys}{3}
              c_1  &-  &c_2  &+  &c_3  &=  &2  \\
              c_1  &+  &c_2  &-  &c_3  &=  &0  \\
             -c_1  &-  &c_2  &+  &c_3  &=  &0
            \end{linsys}
            \grstep[\rho_1+\rho_3]{-\rho_1+\rho_2}
            \begin{linsys}{3}
              c_1  &-  &c_2  &+  &c_3  &=  &2  \\
                   &   &2c_2 &-  &2c_3 &=  &-2  \\
                   &   &-2c_2&+  &2c_3 &=  &2
            \end{linsys}
            \grstep{\rho_2+\rho_3}
            \begin{linsys}{3}
              c_1  &-  &c_2  &+  &c_3  &=  &2  \\
                   &   &2c_2 &-  &2c_3 &=  &-2  \\
                   &   &     &   &0    &=  &0
            \end{linsys}
          \end{equation*}
          给出了额外的信息（除了解至少存在一个之外）：有无穷多解。
          参数化给出 $c_2=-1+c_3$ 与 $c_1=1$，取 $c_3$ 为
          零即得一个特解 $c_1=1$，$c_2=-1$，
          $c_3=0$（这当然就是开头所作的观察）。
      \end{exparts}
    \end{answer}
  \recommended \item
    求此矩阵行空间的一个基。
    \begin{equation*}
      \begin{mat}[r]
         2  &0  &3  &4  \\
         0  &1  &1  &-1 \\
         3  &1  &0  &2  \\
         1  &0  &-4 &-1
       \end{mat}
    \end{equation*}
    \begin{answer}
     % Using Sage:
     % M = matrix(QQ, [[2,0,3,4], [0,1,1,-1], [3,1,0,2], [1,0,-4,-1]])
     % M1 = M.with_added_multiple_of_row(2,0,-3/2)
     % M2 = M1.with_added_multiple_of_row(3,0,-1/2)
     % M3 = M2.with_added_multiple_of_row(2,1,-1)
     % M4 = M3.with_added_multiple_of_row(3,2,-1)
     常规的高斯消去
     \begin{equation*}
       \begin{mat}[r]
         2  &0  &3 &4   \\
         0  &1  &1  &-1 \\
         3  &1  &0  &2  \\
         1  &0  &-4  &-1
       \end{mat}
       \grstep[-(1/2)\rho_1+\rho_4]{-(3/2)\rho_1+\rho_3}
       \grstep{-\rho_2+\rho_3}
       \repeatedgrstep{-\rho_3+\rho_4}
       \begin{mat}[r]
         2  &0  &3      &4   \\
         0  &1  &1      &-1 \\
         0  &0  &-11/2  &-3  \\
         0  &0  &0      &0
       \end{mat}
     \end{equation*}
     提示了这个基
     $\sequence{\rowvec{2 &0 &3 &4},
         \rowvec{0 &1 &1 &-1},\rowvec{0 &0 &-11/2 &-3}}$。

      另一个或许更方便的做法是先对换两行，
      % From Sage
      % M = matrix(QQ, [[2,0,3,4], [0,1,1,-1], [3,1,0,2], [1,0,-4,-1]])
      % M1 = M.with_swapped_rows(0,3)
      % M2 = M1.with_added_multiple_of_row(2,0,-3)
      % M3 = M2.with_added_multiple_of_row(3,0,-2)
      % M4 = M3.with_added_multiple_of_row(2,1,-1)
      % M5 = M4.with_added_multiple_of_row(3,2,-1)
      \begin{equation*}
        \grstep{\rho_1\swap\rho_4}
        \repeatedgrstep[-2\rho_1+\rho_4]{-3\rho_1+\rho_3}
        \repeatedgrstep{-\rho_2+\rho_3}
        \repeatedgrstep{-\rho_3+\rho_4}
        \begin{mat}[r]
           1  &0  &-4 &-1 \\
           0  &1  &1  &-1 \\
           0  &0  &11 &6  \\
           0  &0  &0  &0
         \end{mat}
      \end{equation*}
      从而得到基
      \( \sequence{\rowvec{1 &0 &-4 &-1},\rowvec{0 &1 &1 &-1},
                   \rowvec{0 &0 &11 &6}} \)。
    \end{answer}
  \recommended \item
    求每个矩阵的秩。
    \begin{exparts*}
      \partsitem \(
        \begin{mat}[r]
          2  &1  &3  \\
          1  &-1 &2  \\
          1  &0  &3
        \end{mat}  \)
      \partsitem \(
        \begin{mat}[r]
          1  &-1 &2  \\
          3  &-3 &6  \\
         -2  &2  &-4
        \end{mat}  \)
      \partsitem \(
        \begin{mat}[r]
          1  &3  &2  \\
          5  &1  &1  \\
          6  &4  &3
        \end{mat}  \)
      \partsitem \(
        \begin{mat}[r]
          0  &0  &0  \\
          0  &0  &0  \\
          0  &0  &0
        \end{mat}  \)
    \end{exparts*}
    \begin{answer}
      \begin{exparts}
         \partsitem 这一化简
           \begin{equation*}
             \mbox{}
             \grstep[-(1/2)\rho_1+\rho_3]{-(1/2)\rho_1+\rho_2}
             \repeatedgrstep{-(1/3)\rho_2+\rho_3}
             \begin{mat}[r]
               2  &1     &3     \\
               0  &-3/2  &1/2   \\
               0  &0     &4/3
             \end{mat}
           \end{equation*}
           表明行秩（从而秩）为三。
         \partsitem 观察各列可见其余各列都是第一列的倍数（观察各行也可看出同样的结论）。
           因此秩为一。

           或者，化简
           \begin{equation*}
             \begin{mat}[r]
               1  &-1  &2  \\
               3  &-3  &6  \\
               -2 &2   &-4
             \end{mat}
             \grstep[2\rho_1+\rho_3]{-3\rho_1+\rho_2}
             \begin{mat}[r]
               1  &-1  &2  \\
               0  &0   &0  \\
               0  &0   &0
             \end{mat}
           \end{equation*}
           也表明同样的事实。
         \partsitem 这一计算
           \begin{equation*}
             \begin{mat}[r]
               1  &3  &2  \\
               5  &1  &1  \\
               6  &4  &3
             \end{mat}
             \grstep[-6\rho_1+\rho_3]{-5\rho_1+\rho_2}
             \repeatedgrstep{-\rho_2+\rho_3}
             \begin{mat}[r]
               1  &3   &2  \\
               0  &-14 &-9 \\
               0  &0   &0
             \end{mat}
           \end{equation*}
           表明秩为二。
         \partsitem 秩为零。
       \end{exparts}
     \end{answer}
  \item 给出此矩阵列空间的一个基。
    给出该矩阵的秩。
    \begin{equation*}
      \begin{mat}
        1 &3 &-1 &2 \\
        2 &1 &1  &0 \\
        0 &1 &1  &4
      \end{mat}
    \end{equation*}
  \begin{answer}
    我们要求这个张成空间的一个基。
    \begin{equation*}
      \spanof{\colvec{1 \\ 2 \\ 0},
              \colvec{3 \\ 1 \\ 1},
              \colvec{-1 \\ 1 \\ 1},
              \colvec{2 \\ 0 \\ 4}}\subseteq\Re^3
    \end{equation*}
    最直接的做法
    是把这些列转置成行，用高斯消元法求出这些行的张成空间的一个基，
    然后再把它们转置回
    列。
    \begin{equation*}
      \begin{mat}
        1 &2 &0 \\
        3 &1 &1 \\
       -1 &1 &1 \\
        2 &0 &4
      \end{mat}
      \grstep[\rho_1+\rho_3 \\ -2\rho_1+\rho_4]{-3\rho_1+\rho_2}
      \grstep[-(4/5)\rho_2+\rho_4]{(3/5)\rho_2+\rho_3}
      \grstep{-2\rho_3+\rho_4}
      \begin{mat}
        1 &2  &0 \\
        0 &-5 &1 \\
        0 &0  &8/5 \\
        0 &0  &0
      \end{mat}
    \end{equation*}
    舍去零向量（它表明初始向量之间存在冗余），得到列空间的这个基。
    \begin{equation*}
      \sequence{
        \colvec{1 \\ 2 \\ 0},
        \colvec{0 \\ -5 \\ 1},
        \colvec{0 \\ 0 \\ 8/5}
        }
    \end{equation*}
    矩阵的秩是其列空间的维数，所以为三。
    （它也等于其行空间的维数。）
  \end{answer}
  \recommended \item
    求以下每个集合的张成空间的一个基。
    \begin{exparts}
      \partsitem \(
        \set{\rowvec{1 &3},
          \rowvec{-1 &3},
          \rowvec{1 &4},
          \rowvec{2 &1}  }\subseteq\matspace_{\nbym{1}{2}} \)
      \partsitem \(
         \set{\colvec[r]{1 \\2 \\1},
           \colvec[r]{3 \\ 1 \\ -1},
           \colvec[r]{1 \\ -3 \\ -3}  }\subseteq\Re^3  \)
      \partsitem \(  \set{1+x,1-x^2,3+2x-x^2}\subseteq\polyspace_3  \)
      \partsitem \( \set{
          \begin{mat}[r]
            1  &0  &1  \\
            3  &1  &-1
          \end{mat},
          \begin{mat}[r]
            1  &0  &3  \\
            2  &1  &4
          \end{mat},
          \begin{mat}[r]
           -1  &0  &-5 \\
           -1  &-1 &-9
          \end{mat}  }  \subseteq\matspace_{\nbym{2}{3}}  \)
    \end{exparts}
    \begin{answer}
      \begin{exparts}
        \partsitem 这一化简
          \begin{equation*}
             \begin{mat}[r]
               1  &3  \\
              -1  &3  \\
               1  &4  \\
               2  &1
             \end{mat}
             \grstep[-\rho_1+\rho_3 \\ -2\rho_1+\rho_4]{\rho_1+\rho_2}
             \grstep[(5/6)\rho_2+\rho_4]{-(1/6)\rho_2+\rho_3}
             \begin{mat}[r]
               1  &3  \\
               0  &6  \\
               0  &0  \\
               0  &0
             \end{mat}
          \end{equation*}
          给出 \( \sequence{\rowvec{1 &3},\rowvec{0 &6}} \)。
        \partsitem 转置并化简
          \begin{equation*}
             \begin{mat}[r]
               1  &2  &1  \\
               3  &1  &-1 \\
               1  &-3 &-3
             \end{mat}
             \grstep[-\rho_1+\rho_3]{-3\rho_1+\rho_2}
             \begin{mat}[r]
               1  &2  &1  \\
               0  &-5 &-4 \\
               0  &-5 &-4
             \end{mat}
             \grstep{-\rho_2+\rho_3}
             \begin{mat}[r]
               1  &2  &1  \\
               0  &-5 &-4 \\
               0  &0  &0
             \end{mat}
          \end{equation*}
          再转置回去，得到这个基。
          \begin{equation*}
            \sequence{\colvec[r]{1 \\ 2 \\ 1},
              \colvec[r]{0 \\ -5 \\ -4}   }
          \end{equation*}
        \partsitem 首先注意外围空间是
          $\polyspace_3$，而不是 $\polyspace_2$。
          然后，把第一个多项式 $1+1\cdot x+0\cdot x^2+0\cdot x^3$
          取作与行向量 $\rowvec{1 &1 &0 &0}$“相同”，等等，
          就得到
          \begin{equation*}
            \begin{mat}[r]
              1  &1  &0  &0 \\
              1  &0  &-1 &0 \\
              3  &2  &-1 &0
            \end{mat}
            \grstep[-3\rho_1+\rho_3]{-\rho_1+\rho_2}
            \repeatedgrstep{-\rho_2+\rho_3}
            \begin{mat}[r]
              1  &1  &0  &0 \\
              0  &-1 &-1 &0 \\
              0  &0  &0  &0
            \end{mat}
          \end{equation*}
          由此得到基 \( \sequence{1+x,-x-x^2} \)。
        \partsitem 这里“相同”的做法给出
          \begin{equation*}
            \begin{mat}[r]
              1  &0  &1  &3  &1  &-1  \\
              1  &0  &3  &2  &1  &4   \\
             -1  &0  &-5  &-1 &-1 &-9
            \end{mat}
            \grstep[\rho_1+\rho_3]{-\rho_1+\rho_2}
            \repeatedgrstep{2\rho_2+\rho_3}
            \begin{mat}[r]
              1  &0  &1  &3  &1  &-1  \\
              0  &0  &2  &-1 &0  &5   \\
              0  &0  &0  &0  &0  &0
            \end{mat}
          \end{equation*}
          从而得到这个基。
          \begin{equation*}
            \sequence{
            \begin{mat}[r]
              1  &0  &1  \\
              3  &1  &-1
            \end{mat},
            \begin{mat}[r]
              0  &0  &2  \\
             -1  &0  &5
            \end{mat}  }
          \end{equation*}
      \end{exparts}
     \end{answer}
  \item 在自然的向量空间中，给出以下每个集合的张成空间的一个基。
    \begin{exparts}
      \partsitem
         $\set{\colvec{1 \\ 1 \\ 3},
              \colvec{-1 \\ 2 \\ 0},
              \colvec{0  \\ 12 \\ 6}}$
      \partsitem $\set{x+x^2, 2-2x, 7, 4+3x+2x^2}$
    \end{exparts}
    \begin{answer}
      \begin{exparts}
        \partsitem
          把这些列转置成行，化为阶梯形
          （然后去掉零行），再转置回列。
          \begin{equation*}
            \begin{mat}
              1 &1  &3 \\
             -1 &2  &0 \\
              0 &12 &6
            \end{mat}
            \grstep{\rho_1+\rho_2}
            \grstep{-4\rho_2+\rho_3}
            \begin{mat}
              1 &1  &3 \\
              0 &3  &3 \\
              0 &0  &-6
            \end{mat}
          \end{equation*}
          该张成空间的一个基如下。
          \begin{equation*}
            \sequence{
              \colvec{1 \\ 1 \\ 3},
              \colvec{0 \\ 3 \\ 3},
              \colvec{0 \\ 0 \\ -6}
          }
          \end{equation*}
       \partsitem
          与前一小区一样，我们把它们看作行，
          以利用高斯消元法已完成的工作。
          \begin{multline*}
            \begin{mat}
              0 &1  &1 \\
              2 &-2 &0 \\
              7 &0 &0 \\
              4 &3  &2
            \end{mat}
            \grstep{\rho_1\leftrightarrow\rho_2}
            \grstep[2\rho_1+\rho_4]{-(7/2)\rho_1+\rho_3}
            \grstep[-7\rho_2+\rho_4]{-7\rho_2+\rho_3}          \\
            \grstep{-(5/7)\rho_3+\rho_4}
            \begin{mat}
              2 &-2  &0 \\
              0 &1 &1 \\
              0 &0  &-7 \\
              0 &0 &0
            \end{mat}
          \end{multline*}
          该集合的张成空间的一个基是
          $\sequence{2-2x, x+x^2, -5x^2}$。
      \end{exparts}
    \end{answer}
  \item
    哪些矩阵的秩为零？
    秩为一的呢？
    \begin{answer}
      只有零矩阵的秩为零。
      秩为一的矩阵仅有如下形式
      \begin{equation*}
        \begin{mat}
           k_1\cdot \rho  \\
            \vdots  \\
           k_m\cdot \rho
        \end{mat}
      \end{equation*}
      其中 \( \rho \) 是某个非零行向量，且诸 \( k_i \) 不全为零。
      （\textit{注}：
      我们不能简单地说所有行都是第一行的倍数，因为第一行可能是零行。
      \textit{再注}：
      上述结论在把“行”换成“列”后同样适用。）
    \end{answer}
  \recommended \item
    已知 \( a,b,c\in\Re \)，要取什么样的 \( d \) 才能使这个矩阵
    的秩为一？
    \begin{equation*}
      \begin{mat}
         a  &b  \\
         c  &d
      \end{mat}
    \end{equation*}
    \begin{answer}
      若 \( a\neq 0 \)，则取 \( d=(c/a)b \) 可使第二行是第一行的倍数，
      具体地说，是第一行的 \( c/a \) 倍。
      若 \( a=0 \) 且 \( b=0 \)，则 \( d \) 取任意非 \( 0 \) 值都能保证
      第二行非零。
      若 \( a=0 \) 且 \( b\neq 0 \) 且 \( c=0 \)，则 \( d \) 取任何值都行，
      因为该矩阵自动具有秩一（即使取 $d=0$）。
      最后，若 \( a=0 \) 且 \( b\neq 0 \) 且 \( c\neq 0 \)，则 \( d \) 取
      任何值都不够，因为该矩阵必定有秩二。
    \end{answer}
  \item
    求这个矩阵的列秩。
    \begin{equation*}
      \begin{mat}[r]
        1  &3  &-1  &5  &0  &4  \\
        2  &0  &1   &0  &4  &1
      \end{mat}
    \end{equation*}
    \begin{answer}
      列秩是二。
      一种看法是直接观察\Dash 列空间由高度为二的列组成，因而其维数
      至少可以是二，而且我们容易找到两列，它们合起来组成一个
      线性无关的集合（例如第四列与第五列）。
      另一种看法是回忆列秩等于行秩，然后施行高斯消元法，
      结果剩下两个非零行。
    \end{answer}
  \item
    证明：一个至少有一个解的线性方程组至多有一个解当且仅当
    其系数矩阵的秩等于其列数。
    \begin{answer}
      我们应用\nearbytheorem{th:RankVsSoltnSp}。
      线性方程组的系数矩阵 $A$ 的列数等于未知数的个数 $n$。
      一个至少有一个解的线性方程组至多有一个解当且仅当相应的
      齐次方程组的解空间维数为零（回忆：~在
      “$\text{通解}=\text{特解}+\text{齐次解}$”这一方程
       $\vec{v}=\vec{p}+\vec{h}$ 中，只要这样的 $\vec{p}$ 存在，
       解 $\vec{v}$ 唯一当且仅当向量 $\vec{h}$ 唯一，即 $\vec{h}=\zero$）。
       而由该定理，这意味着 $n=r$。
    \end{answer}
  \recommended \item
    如果一个矩阵是 \( \nbym{5}{9} \) 的，它的行集合与列集合当中，
    哪一个必定是线性相关的？
    \begin{answer}
      列的集合必定是线性相关的，因为该矩阵的秩至多是五，而列却有
      九个。
    \end{answer}
  \item
    举一个例子，说明尽管两者维数相同，一个矩阵的行空间与列空间
    却不必相等。
    它们有可能相等吗？
    \begin{answer}
      它们几乎没有相等的可能，因为行空间是行向量的集合，而列空间是
      列向量的集合（除非该矩阵是 \( \nbyn{1} \) 的，此时两个空间
      必定相等）。

      \textit{注记。}
      考察
      \begin{equation*}
        A=\begin{mat}[r]
            1  &3  \\
            2  &6
          \end{mat}
      \end{equation*}
      并注意行空间是 \( \rowvec{1 &3} \) 的所有倍数组成的集合，
      而列空间由以下列向量的倍数组成
      \begin{equation*}
        \colvec[r]{1 \\ 2}
      \end{equation*}
      所以我们也不能论证说这两个空间必定恰好互为转置。
    \end{answer}
  \item
    证明集合
    \( \set{(1,-1,2,-3),(1,1,2,0),(3,-1,6,-6)} \) 与
    \( \set{(1,0,1,0),(0,2,0,3)} \) 的张成不相同。
    顺便问一句，这里的向量空间是什么？
    \begin{answer}
      首先，该向量空间是在自然运算下的实数四元组的集合。
      尽管这不是四宽的行向量的集合，但差别不大\Dash 它与那个集合
      是“相同的”。
      所以我们把四元组当作四宽的向量来处理。

      有了这一点，要看出 \( (1,0,1,0) \) 不在第一个集合的张成中，
      一种办法是注意这个化简
      \begin{equation*}
        \begin{mat}[r]
          1  &-1  &2  &-3  \\
          1  &1   &2  &0   \\
          3  &-1  &6  &-6
        \end{mat}
        \grstep[-3\rho_1+\rho_3]{-\rho_1+\rho_2}
        \repeatedgrstep{-\rho_2+\rho_3}
        \begin{mat}[r]
          1  &-1  &2  &-3  \\
          0  &2   &0  &3   \\
          0  &0   &0  &0
        \end{mat}
      \end{equation*}
      以及这个化简
      \begin{equation*}
        \begin{mat}[r]
          1  &-1  &2  &-3  \\
          1  &1   &2  &0   \\
          3  &-1  &6  &-6  \\
          1  &0   &1  &0
        \end{mat}
        \grstep[-3\rho_1+\rho_3 \\ -\rho_1+\rho_4]{-\rho_1+\rho_2}
        \repeatedgrstep[-(1/2)\rho_2+\rho_4]{-\rho_2+\rho_3}
        \repeatedgrstep{\rho_3\leftrightarrow\rho_4}
        \begin{mat}[r]
          1  &-1  &2  &-3  \\
          0  &2   &0  &3   \\
          0  &0   &-1 &3/2 \\
          0  &0   &0  &0
        \end{mat}
      \end{equation*}
      得到的矩阵在秩上不同。
      这意味着把 $(1,0,1,0)$ 添加到前三个四元组的集合中会增大该
      集合的秩，从而增大其张成。
      因此 $(1,0,1,0)$ 原本就不在这个张成中。
    \end{answer}
  \recommended \item
    证明这个列向量的集合
    \begin{equation*}
      \set{\colvec{d_1 \\ d_2 \\ d_3}
           \suchthat
           \text{存在 $x$、$y$ 与 $z$ 使得：\ }
           \begin{linsys}{3}
             3x  &+  &2y  &+  &4z  &=   &d_1   \\
              x  &   &    &-  &z   &=   &d_2   \\
             2x  &+  &2y  &+  &5z  &=   &d_3
           \end{linsys}
           }
    \end{equation*}
    是 \( \Re^3 \) 的一个子空间。
    求一个基。
    \begin{answer}
      它是子空间，因为它是系数矩阵
      \begin{equation*}
        \begin{mat}[r]
          3  &2  &4  \\
          1  &0  &-1 \\
          2  &2  &5
        \end{mat}
      \end{equation*}
      的列空间。
      为了求列空间的一个基，
      \begin{equation*}
        \set{c_1\colvec[r]{3 \\ 1 \\ 2}
             +c_2\colvec[r]{2 \\ 0 \\ 2}
             +c_3\colvec[r]{4 \\ -1 \\ 5}
             \suchthat c_1,c_2,c_3\in\Re}
      \end{equation*}
      我们通过转置、化简
      \begin{equation*}
         \begin{mat}[r]
           3  &1  &2  \\
           2  &0  &2  \\
           4  &-1 &5
         \end{mat}
         \grstep[-(4/3)\rho_1+\rho_3]{-(2/3)\rho_1+\rho_2}
         \repeatedgrstep{-(7/2)\rho_2+\rho_3}
         \begin{mat}[r]
           3  &1     &2  \\
           0  &-2/3  &2/3  \\
           0  &0     &0
         \end{mat}
      \end{equation*}
      并略去零行，再转置回去，来消去张成集中三个列向量之间的
      线性关系。
      \begin{equation*}
        \sequence{ \colvec[r]{3 \\ 1 \\ 2},
                   \colvec[r]{0 \\ -2/3 \\ 2/3}  }
      \end{equation*}
    \end{answer}
  \item
    证明转置运算是
    线性的：\index{transpose!is linear}
    \begin{equation*}
      \trans{(rA+sB)}  = r\trans{A}+s\trans{B}
    \end{equation*}
    其中 \( r,s\in\Re \) 且 \( A,B\in\matspace_{\nbym{m}{n}} \)。
    \begin{answer}
      我们可以通过直接的计算来完成。
      \begin{align*}
        \trans{(rA+sB)}
        &=\trans{\begin{mat}
             ra_{1,1}+sb_{1,1}  &\ldots &ra_{1,n}+sb_{1,n} \\
                          &\vdots            \\
             ra_{m,1}+sb_{m,1}  &\ldots &ra_{m,n}+sb_{m,n}
                \end{mat}  }  \\
        &=\begin{mat}
           ra_{1,1}+sb_{1,1}  &\ldots &ra_{m,1}+sb_{m,1}  \\
                            &\vdots                           \\
           ra_{1,n}+sb_{1,n}  &\ldots &ra_{m,n}+sb_{m,n}
         \end{mat}           \\
        &=\begin{mat}
           ra_{1,1}  &\ldots &ra_{m,1}  \\
                    &\vdots                       \\
           ra_{1,n}  &\ldots &ra_{m,n}
         \end{mat}
        +\begin{mat}
           sb_{1,1}  &\ldots &sb_{m,1}  \\
                    &\vdots                       \\
           sb_{1,n}  &\ldots &sb_{m,n}
         \end{mat}           \\
        &=r\trans{A}+s\trans{B}
      \end{align*}
    \end{answer}
  \recommended \item
    在本小节中我们已经证明高斯化简能求出行空间的一个基。
    \begin{exparts}
      \partsitem 证明这个基不是唯一的\Dash 不同的化简可能给出
        不同的基。
      \partsitem 构造行空间相等但行数不相等的矩阵。
      \partsitem 证明两个矩阵有相同的行空间当且仅当经过
        高斯–若尔当消元后它们有相同的非零行。
    \end{exparts}
    \begin{answer}
      \begin{exparts}
        \partsitem 这些化简给出不同的基。
          \begin{equation*}
            \begin{mat}[r]
              1  &2  &0  \\
              1  &2  &1
            \end{mat}
            \grstep{-\rho_1+\rho_2}
            \begin{mat}[r]
              1  &2  &0  \\
              0  &0  &1
            \end{mat}
            \qquad
            \begin{mat}[r]
              1  &2  &0  \\
              1  &2  &1
            \end{mat}
            \grstep{-\rho_1+\rho_2}
            \repeatedgrstep{2\rho_2}
            \begin{mat}[r]
              1  &2  &0  \\
              0  &0  &2
            \end{mat}
          \end{equation*}
        \partsitem 一个简单的例子是这个。
          \begin{equation*}
            \begin{mat}[r]
              1  &2  &1  \\
              3  &1  &4
            \end{mat}
            \qquad
            \begin{mat}[r]
              1  &2  &1  \\
              3  &1  &4  \\
              0  &0  &0
            \end{mat}
          \end{equation*}
          这是一个稍微不那么简单的例子。
          \begin{equation*}
            \begin{mat}[r]
              1  &2  &1  \\
              3  &1  &4
            \end{mat}
            \qquad
            \begin{mat}[r]
              1  &2  &1  \\
              3  &1  &4  \\
              2  &4  &2  \\
              4  &3  &5
            \end{mat}
          \end{equation*}
        \partsitem
          % We give two proofs.
          % The first is a bit lower-level but more constructive.
          % The second is higher level.

          % \textit{First proof.}
          % Assume that \( A \) and \( B \) are matrices with equal row spaces.
          % Construct a matrix \( C \) with the rows of \( A \) above the rows
          % of \( B \), and another matrix \( D \) with the rows of \( B \)
          % above the rows of \( A \).
          % \begin{equation*}
          %   C=\begin{mat}
          %        A \\ B
          %     \end{mat}
          %   \qquad
          %   D=\begin{mat}
          %        B \\ A
          %     \end{mat}
          % \end{equation*}
          % Observe that \( C \) and \( D \) are row-equivalent (via a sequence
          % of row-swaps) and so Gauss-Jordan reduce to the same reduced
          % echelon form matrix.

          % Because the row spaces are equal, the rows of \( B \)
          % are linear combinations of the rows of \( A \) so Gauss-Jordan
          % reduction on \( C \) simply turns the rows of \( B \) to zero rows
          % and thus the nonzero rows of $C$ are just the nonzero rows obtained
          % by Gauss-Jordan reducing \( A \).
          % The same can be said for the matrix \( D \)\Dash Gauss-Jordan
          % reduction on \( D \) gives the same non-zero rows as are produced
          % by reduction on \( B \) alone.
          % Therefore,
          % \( A \) yields the same nonzero rows as \( C  \), which yields
          % the same nonzero rows as \( D \), which yields the same nonzero
          % rows as \( B \).

          % \textit{Second proof (N~Schenker).}
          由于 \( A \) 与 \( B \) 的行空间相等，所以二者行等价。
          由于每个行等价类都包含唯一的简化阶梯形
          （定理 One.III.\ref{th:ReducedEchelonFormIsUnique}），
          \( A \) 的简化阶梯形必定等于 \( B \) 的简化阶梯形。
      \end{exparts}
    \end{answer}
  \item
    对于 \nearbyremark{rem:MunkresCommment}，在 \( r \) 大于 \( n \) 的
    情形下为什么不成问题？
    \begin{answer}
      它不可能更大。
    \end{answer}
  \item
    证明 \( \nbym{m}{n} \) 矩阵的行秩至多是 \( m \)。
    有更好的界吗？
    \begin{answer}
      极大线性无关集合中行的个数不会超过行的个数。
      一个更好的界（一般情形下这是最佳可能的界）是 \( m \) 与 \( n \)
      中的最小者，因为行秩等于列秩。
     \end{answer}
  \item
    证明一个矩阵的秩等于其转置的秩。
    \begin{answer}
      由于矩阵 \( A \) 的行是 \( \trans{A} \) 的列，所以 \( A \) 的
      行空间的维数等于 \( \trans{A} \) 的列空间的维数。
      而 \( A \) 的行空间的维数就是 \( A \) 的秩，并且 \( \trans{A} \) 的
      列空间的维数就是 \( \trans{A} \) 的秩。
      于是这两个秩相等。
    \end{answer}
  \item
    判断真假：~一个矩阵的列空间等于其转置的行空间。
    \begin{answer}
      假的。
      前者是列的集合，而后者是行的集合。

      然而这个例子
      \begin{equation*}
         A=\begin{mat}[r]
             1  &2  &3  \\
             4  &5  &6
           \end{mat},
         \qquad
         \trans{A}=\begin{mat}[r]
             1  &4  \\
             2  &5  \\
             3  &6
           \end{mat}
      \end{equation*}
      表明，一旦我们对“相同”有了形式化的含义，就可以在这里应用它：
      \begin{equation*}
        \colspace{A}=\spanof{\set{
                        \colvec[r]{1 \\ 4},
                        \colvec[r]{2 \\ 5},
                        \colvec[r]{3 \\ 6} }  }
      \end{equation*}
      而
      \begin{equation*}
         \rowspace{\trans{A}}=\spanof{\set{
                                 \rowvec{1 &4},
                                 \rowvec{2 &5},
                                 \rowvec{3 &6} }  }
      \end{equation*}
      彼此是“相同的”。
    \end{answer}
  \recommended \item
    我们已经看到行变换可能改变列空间。
    它必定会吗？
    \begin{answer}
      不会。
      这里，高斯消元法没有改变列空间。
      \begin{equation*}
        \begin{mat}[r]
          1  &0  \\
          3  &1
        \end{mat}
        \grstep{-3\rho_1+\rho_2}
        \begin{mat}[r]
          1  &0  \\
          0  &1
        \end{mat}
      \end{equation*}
    \end{answer}
  \item
    证明一个线性方程组有解当且仅当该方程组的系数矩阵与
    其增广矩阵有相同的秩。
    \begin{answer}
      线性方程组
      \begin{equation*}
        c_1\vec{a}_1+\dots+c_n\vec{a}_n=\vec{d}
      \end{equation*}
      有解当且仅当 \( \vec{d} \) 在集合 \( \set{\vec{a}_1,\dots,\vec{a}_n} \)
      的张成中。
      这成立当且仅当增广矩阵的列秩等于系数矩阵的列秩。
      由于秩等于列秩，方程组有解当且仅当其增广矩阵的秩等于其
      系数矩阵的秩。
    \end{answer}
  \item
    一个 \( \nbym{m}{n} \) 矩阵具有
    \definend{行满秩}\index{full row rank}\index{row rank!full}，
    若其行秩
    为 \( m \)；它具有
    \definend{列满秩}\index{full column rank}\index{column!rank!full}，
    若其列秩为 \( n \)。
    \begin{exparts}
      \partsitem 证明一个矩阵能同时有行满秩与列满秩仅当它是方阵。
      \partsitem 证明：系数矩阵为 \( A \) 的线性方程组对右侧任意的
        \( d_1 \)、\ldots、\( d_n \) 都有解当且仅当 \( A \) 行满秩。
      \partsitem 证明齐次方程组有唯一解当且仅当其系数矩阵
        $A$ 列满秩。
      \partsitem 证明命题
       “若系数矩阵为 \( A \) 的方程组有任何解则它有唯一解”成立
        当且仅当 \( A \) 列满秩。
    \end{exparts}
    \begin{answer}
      \begin{exparts}
        \partsitem 行秩等于列秩，所以二者各自至多是行数与列数的
          最小者。
          因此只有当行数等于列数时二者才能都满。
          （当然，逆命题不成立：~方阵不必行满秩或列满秩。）
        \partsitem 若 \( A \) 行满秩，则无论右侧是什么，对增广矩阵施行
          高斯消元法结束时每一行都有一个首主元一，并且这些首主元一
          没有一个位于最右边的列（“增广”列）。
          回代随即给出一个解。

          另一方面，如果该线性方程组对某个右侧没有解，只可能是因为
          高斯消元法留下某行，使得它在“增广”竖线左边全为零，
          而在右边有非零元。
          于是，如果 $A$ 对某些右侧没有解，那么 \( A \) 不具有
          行满秩，因为它的某些行已被消去。
        \partsitem 矩阵 \( A \) 列满秩当且
          仅当其列
          组成一个线性无关的集合。
          这等价于诸列之间只存在平凡的线性关系，于是方程组唯一的解是
          每个变量都取 $0$ 的情形。
        \partsitem 矩阵 \( A \) 列满秩当且
          仅当其列的集合是线性无关的，从而构成其张成的一个基。
          这等价于该张成中的所有向量都有唯一的线性表示。
          这就证明了它，
          因为张成中一个向量的任何线性表示都是
          该线性方程组的解。
      \end{exparts}
    \end{answer}
  \item
    如果把高斯消元法改成允许用零去乘某一行，
    \nearbylemma{le:RowSpUnchByGR} 的结论会如何改变？
    \begin{answer}
      行空间不再相同，$B$ 的行空间将是 $A$ 的行空间的一个子空间
      （有可能相等）。
    \end{answer}
  \item
    \( \rank(A) \) 与 \( \rank(-A) \) 之间有什么关系？
    \( \rank(A) \) 与 \( \rank(kA) \) 之间呢？
    \( \rank(A) \)、\( \rank(B) \) 与 \( \rank(A+B) \) 之间，
    如果有关系的话，是什么关系？
    \begin{answer}
      显然 \( \rank(A)=\rank(-A) \)，因为高斯消元法允许我们把一个
      矩阵的所有行乘以 \( -1 \)。
      同样地，当 \( k\neq 0 \) 时我们有 \( \rank(A)=\rank(kA) \)。

      加法更有意思。
      和的秩可以比各加数的秩小。
      \begin{equation*}
        \begin{mat}[r]
          1  &2  \\
          3  &4
        \end{mat}
        +
        \begin{mat}[r]
         -1  &-2  \\
         -3  &-4
        \end{mat}
        =
        \begin{mat}[r]
         0   &0   \\
         0   &0
        \end{mat}
      \end{equation*}
      和的秩也可以比各加数的秩大。
      \begin{equation*}
        \begin{mat}[r]
          1  &2  \\
          0  &0
        \end{mat}
        +
        \begin{mat}[r]
          0  &0   \\
          3  &4
        \end{mat}
        =
        \begin{mat}[r]
         1   &2   \\
         3   &4
        \end{mat}
      \end{equation*}
      但存在一个上界（除矩阵尺寸之外的上界）。
      一般地，\( \rank(A+B)\leq\rank(A)+\rank(B) \)。

      为了证明它，注意我们可以用两种方式之一对 \( A+B \) 实施
      高斯消去法：
      可以先把 \( A \) 加到 \( B \) 上，然后施行适当的化简步骤序列
      \begin{equation*}
        (A+B)
        \grstep{\text{步骤}_1}
        \cdots
        \grstep{\text{步骤}_k}
        \text{阶梯形}
      \end{equation*}
      或者也可以分别在 \( A \) 与 \( B \) 上实施 \( \text{步骤}_1 \) 到
      \( \text{步骤}_k \)，然后再相加，得到同样的结果。
      在第二种方式中我们最终能得到的最大秩显然是两个秩之和。
      （上面的矩阵给出了两种可能性，
       $\rank(A+B)<\rank(A)+\rank(B)$ 与 $\rank(A+B)=\rank(A)+\rank(B)$，
       都发生的例子。）
    \end{answer}
    % Relationship between rank(A), rank(B) and rank(cA+dB)?
\end{exercises}























